% MEJORADO (27 oct 2026): realidad aumentada acotada (Azuma, 1997), umbral y precisión del
% reconocimiento, componente de servidor, límites de las pruebas y trabajo futuro con fuentes.
\chapter{CONCLUSIONES Y RECOMENDACIONES}
\label{cap:conclusiones}

\section{Conclusiones}
\label{sec:conclusiones}

El presente proyecto de grado desarrolló un sistema para el entrenamiento en
gimnasios que proporciona una guía estructurada y un seguimiento personalizado
de los ejercicios, con lo que se cumple el objetivo general planteado. El
producto es una aplicación Android construida con React Native, Expo y
TypeScript sobre los servicios de Firebase, organizada en una arquitectura
cliente--servidor con capas Modelo--Vista--Controlador y desarrollada con la
metodología Mobile-D en seis iteraciones. Dos de sus componentes, la
inteligencia artificial y la realidad aumentada, se implementaron con un
alcance acotado, como se detalla a continuación: la primera mediante modelos de
TensorFlow Lite entrenados para este proyecto y un detector YOLO26m publicado en Ultralytics Platform, cuya validación se limita a datos sintéticos, de
catálogo y a las particiones del propio entrenamiento; la segunda mediante un módulo
nativo basado en ARCore que ancla la ficha de la máquina en el entorno de los
teléfonos compatibles y, en los demás, una superposición en pantalla sobre la
cámara. Según la definición de Azuma (1997), el anclaje aporta el registro
tridimensional, pero aún no se ha validado en un teléfono físico y el
reconocimiento sigue ocurriendo por captura y no de forma continua
(sección~\ref{subsec:ra-alcance}).

Respecto a cada objetivo específico, se concluye lo siguiente:

\begin{enumerate}[label=\textbf{OE\arabic*.}, leftmargin=3.2em]
    \item \textbf{Guiar sobre el uso correcto de máquinas.} Se construyó un
    catálogo de 35 tipos de máquinas, con 288 modelos comerciales como
    variantes, en el que cada ficha incluye descripción, músculos trabajados,
    guía de uso paso a paso, errores comunes y pautas de prevención de
    lesiones redactadas como contenido propio. El catálogo se consulta con
    búsqueda tolerante a tildes y filtros, y se complementa con un escáner que
    superpone la ficha de la máquina sobre la vista de la cámara y, en los
    teléfonos compatibles con ARCore, la ancla a un punto del entorno. Un
    clasificador de TensorFlow Lite (MobileNetV2 con una capa entrenada),
    entrenado con las 399 fotografías del catálogo, propone la máquina a
    partir de una foto; su acierto es modesto (38,5\,\% en un escenario de
    campo simulado de 35 clases y 61,4\,\% entre las tres primeras
    sugerencias), por lo que solo se acepta sin preguntar con una confianza de
    al menos 0,8, umbral con el que la validación cruzada arrojó una precisión
    del 85,4\,\% sobre el 15,8\,\% de las fotografías; en los demás casos, el
    usuario confirma la sugerencia o elige la máquina en una lista. Como
    alternativa opcional, un detector YOLO26m entrenado y publicado como \emph{endpoint} en Ultralytics Platform devuelve además la caja de la máquina; en validación alcanzó una mAP50 de 48,5\,\% y una exhaustividad de 45,0\,\%, cifras que muestran que el procedimiento funciona de extremo a extremo, no que el detector sea fiable en un gimnasio. El reconocimiento se
    ejecuta sobre una fotografía fija y su latencia en el dispositivo no se ha
    medido.
    \item \textbf{Instruir en la ejecución correcta de ejercicios.} Cada
    máquina presenta al menos un ejercicio con su animación de ejecución, su
    músculo objetivo y su nivel de dificultad, con animaciones y pasos ilustrados procedentes de varias fuentes de datos: WorkoutX (s.~f.), ExerciseDB
    (AscendAPI, s.~f.), Gym Visual (s.~f.), wger (wger Project, s.~f.) y
    free-exercise-db (Yuhonas, s.~f.), que se
    complementan con una guía de entrenamiento paso a paso. El sistema no
    verifica automáticamente la postura del usuario.
    \item \textbf{Diseñar rutinas personalizadas.} Se implementó
    un generador de rutinas basado en reglas cuya base de conocimiento
    traduce las pautas de prescripción del ejercicio en reglas de frecuencia,
    de variables de carga y de división semanal, y que se ajusta con la edad,
    el IMC y los días disponibles del usuario. Sobre esas reglas se entrenó un
    modelo de TensorFlow Lite que decide la plantilla semanal y el esquema de
    carga y se ejecuta en el dispositivo; si su confianza es menor que 0,6 o
    falla, la aplicación usa las reglas. Como no existen datos de usuarios
    reales con resultados conocidos, el modelo se entrenó con 60\,000 perfiles
    sintéticos etiquetados por las propias reglas: reproduce las reglas con un
    99,8\,\% de exactitud, pero eso mide su fidelidad al criterio experto y no
    demuestra que sus rutinas sean mejores para los usuarios. Las pruebas
    sobre las doce combinaciones de objetivo y nivel confirmaron que cada
    rutina generada respeta las reglas.
    \item \textbf{Construir herramientas de planificación.} El planificador
    de rutinas personalizadas permite organizar la semana día por día, ajustar
    series, repeticiones y descansos, mantener varias rutinas con una activa y
    agregar ejercicios directamente desde la ficha de una máquina.
    \item \textbf{Registrar la ingesta de alimentos.} El alcance inicial
    limitaba este objetivo al conteo calórico mediante una calculadora de IMC;
    el sistema lo supera. Además de calcular y clasificar el IMC según la OMS,
    permite buscar alimentos en español sobre la base de USDA FoodData Central,
    calcular calorías y macronutrientes por porción y compararlos con metas
    diarias que se personalizan con la ecuación de Mifflin-St Jeor según los
    datos físicos, los días de entrenamiento y el objetivo del usuario. Son
    estimaciones referenciales y no sustituyen la indicación de un
    nutricionista.
    \item \textbf{Implementar el seguimiento del progreso físico.} El usuario
    registra su peso y su estatura, el sistema calcula el IMC y el historial
    se sincroniza en tiempo real. Además, el usuario marca los días de su
    rutina que completa y el sistema calcula su racha, los días entrenados de
    la semana y el total de entrenamientos.
    \item \textbf{Generar reportes de progreso.} El reporte presenta la
    evolución del peso y del IMC en un gráfico, la variación acumulada y la
    categoría de IMC actual, calculados al vuelo a partir del historial para
    reflejar siempre los datos vigentes. Además, el usuario puede generar un
    reporte en PDF por período (7, 30, 90 días o todo el historial) con su
    perfil, su peso, sus entrenamientos y su alimentación frente a la meta, y
    compartirlo desde el sistema; su generación se verificó en el dispositivo
    de prueba.
\end{enumerate}

En el plano técnico, se concluye además que:

\begin{itemize}
    \item La separación en capas resultó determinante para la mantenibilidad:
    al aislar el acceso a Firebase en los repositorios y servicios, cambios
    como el paso de colecciones raíz a subcolecciones por usuario, o la
    incorporación de nuevas fuentes de datos como USDA, se resolvieron sin
    modificar las pantallas.
    \item Las reglas de seguridad y las funciones del servidor son
    indispensables en una arquitectura basada en servicios gestionados, porque
    las credenciales del cliente son públicas por diseño. La incorporación de
    dos funciones (\archivo{buscarAlimentos} y \archivo{guardarRutina}) con
    App Check, rechazo de usuarios anónimos y cuotas por usuario permitió ocultar
    la clave de USDA y validar en el servidor la creación de rutinas. Esta
    protección aún no cuenta con resultados de ejecución de las pruebas con
    emulador (sección~\ref{sec:pruebas-seguridad}).
    \item Concentrar las reglas de negocio en funciones puras (generador de
    rutinas, IMC, racha de entrenamientos, porciones de alimentos) permitió
    verificarlas de forma automática y reproducible sin depender de la red ni
    de la base de datos.
    \item Aislar los modelos de aprendizaje automático tras un contrato
    definido (las clases que producen y el vector de entrada) y conservar las
    reglas como respaldo permitió integrarlos sin afectar al resto del sistema
    y garantiza que la aplicación siga funcionando si un modelo falla. Además,
    verificar con pruebas automáticas que el orden de las clases en TypeScript
    coincida con el del \emph{script} de entrenamiento evita un error silencioso
    al reentrenar.
    \item Una métrica alta no basta: el recomendador de rutinas alcanza un
    99,8\,\% porque imita reglas conocidas, mientras que el reconocedor de
    máquinas, entrenado con pocas fotografías, acierta menos del 40\,\% en
    campo simulado. Informar ambas cifras con su contexto fue necesario para
    no sobrestimar el aporte de la inteligencia artificial.
\end{itemize}

Como limitaciones del producto actual cabe señalar que las pruebas
automatizadas cubren la lógica pura y los componentes de servicio (20 archivos
de prueba en la aplicación, con 160 casos aprobados el 3 de noviembre de 2026)
pero no la interfaz, que la persistencia de datos sin conexión no está
configurada, que los modelos no se han medido en un dispositivo ni con usuarios
reales, que el anclaje con ARCore no se ha validado en un teléfono físico y que
el ancla no sigue a la máquina si esta se mueve, que los reconocedores aciertan
poco con las máquinas menos distintivas, que el uso del detector en la nube
depende de un servicio externo con licencia propia, y que las metas
nutricionales son estimaciones referenciales.

\section{Recomendaciones}
\label{sec:recomendaciones}

A partir de la experiencia del desarrollo se recomienda, como trabajo futuro:

\begin{itemize}
    \item \textbf{Recomendador de rutinas con datos reales.} Registrar, con
    consentimiento, el perfil de los usuarios, la rutina recibida y su
    adherencia y progreso, para reentrenar el modelo con resultados reales en
    lugar de etiquetas sintéticas y medir si sus rutinas superan a las reglas.
    Las rutinas ya guardan su origen, la versión del modelo y la confianza, lo
    que facilita este análisis.
    \item \textbf{Más fotografías para el reconocimiento de máquinas.} Ampliar
    el conjunto con fotografías propias de las 35 máquinas, tomadas en gimnasios
    reales y con autorización, con varias por máquina (hoy hay desde una hasta
    25), y reentrenar el clasificador y el detector. En el detector, conviene ampliar el conjunto de datos y las clases con menos ejemplos, reservar una partición de prueba independiente y compararlo con otros tamaños de YOLO26 (por ejemplo, el nano, pensado para dispositivos de borde; Jocher et~al., 2026) y con una versión ejecutada en el propio teléfono. Después, recalibrar el umbral de confianza y medir los modelos en dispositivos. No deben usarse las animaciones de
    WorkoutX para el entrenamiento, porque sus condiciones de uso lo prohíben
    sin permiso escrito.
    \item \textbf{Validar y ampliar la realidad aumentada.} El anclaje con
    ARCore ya está implementado, pero falta lo esencial para darlo por
    validado: ejecutar el protocolo de pruebas (CP-42 a CP-44) en teléfonos
    físicos y en un gimnasio, y medir el error de colocación, las pérdidas de
    seguimiento y los modelos de teléfono compatibles. Después se proponen tres
    líneas de trabajo, en orden creciente de complejidad: (a) ejecutar el
    detector sobre los fotogramas del vídeo, no solo sobre una captura, y
    estabilizar el resultado en el tiempo; (b) combinar la caja del detector
    con la profundidad de ARCore en los dispositivos que la admiten, para
    acercar el ancla a la geometría de la máquina; y (c) anclar indicaciones
    tridimensionales (puntos de ajuste, trayectoria del movimiento) sobre la
    máquina real (Google for Developers, s.~f.-f). Revisiones
    del campo (Azuma et al., 2001; Billinghurst et al., 2015) señalan el
    registro preciso y el seguimiento como los problemas técnicos centrales,
    por lo que cada línea debe evaluarse con medidas de latencia y estabilidad
    en un dispositivo físico antes de afirmar que se logró el registro
    tridimensional.
    \item \textbf{Verificación de la postura.} Evaluar la estimación de pose
    corporal en el dispositivo (por ejemplo, BlazePose; Bazarevsky et al.,
    2020) como línea futura para retroalimentar la ejecución de los
    ejercicios, en la línea del antecedente de Silva Amado (2024); queda fuera
    del alcance actual.
    \item \textbf{Progresión de la carga.} Ampliar la base de conocimiento del
    generador con reglas de sobrecarga progresiva que ajusten la rutina según
    el historial de entrenamientos del usuario.
    \item \textbf{Metas nutricionales más precisas.} Permitir que el usuario
    indique su nivel de actividad diaria (hoy se deduce de los días de
    entrenamiento) y ajuste sus metas, manteniendo la advertencia de que no
    sustituyen a un nutricionista.
    \item \textbf{Persistencia sin conexión.} Configurar la caché persistente
    de Firestore para que la guía de máquinas y la rutina activa estén
    disponibles sin Internet, algo muy útil en gimnasios con mala cobertura.
    \item \textbf{Pruebas automatizadas de interfaz y de seguridad.} Agregar
    React Native Testing Library para probar las pantallas, ejecutar la suite de
    Jest y las pruebas con el emulador de Firebase (\archivo{npm run
    test:security}) en integración continua, y registrar sus resultados.
    \item \textbf{Privacidad y consentimiento.} Redactar la política de
    privacidad y el consentimiento informado para el tratamiento de los datos
    físicos y alimentarios, y para el uso de fotografías de usuarios si se
    amplía el conjunto de entrenamiento.

    \item \textbf{Validación con usuarios.} Realizar un estudio de usabilidad
    con usuarios reales de gimnasio, por ejemplo con el cuestionario SUS, y
    medir el efecto de la aplicación en la permanencia de los usuarios nuevos.
    \item \textbf{Licencia de Ultralytics.} La biblioteca y los modelos de
    Ultralytics se distribuyen con licencia AGPL-3.0 o con una licencia
    empresarial (Ultralytics, s.~f.-c). Si la aplicación se publica o se
    comercializa, hay que definir cuál aplica, o sustituir el detector por otra
    alternativa, antes de la distribución.
    \item \textbf{Publicación.} Antes de publicar la aplicación en Google
    Play, obtener la autorización del proveedor de las fotografías del
    catálogo y confirmar las condiciones de uso de las animaciones y los
    datos de ExerciseDB, Gym Visual, WorkoutX y wger, en particular las
    restricciones de uso comercial y de atribución.
\end{itemize}
